\chapter{Architecture}\label{chap:architecture}

\section{Overview}
\soft has two layers:
\begin{description}
  \item[Native core (C++)] format parsing, chip emulation, mixing, per-voice outputs and chip states. It is built
        with the ZXTune \texttt{make} system for four Android architectures into \texttt{libzxtune.so}.
  \item[Application (Kotlin, Java)] playback service, interface, playlists, browsing, oscilloscope rendering (Vulkan
        or OpenGL ES) and dashboards.
\end{description}

\section{Audio path}
\begin{enumerate}
  \item A \emph{Module::Holder} is created for the file; it provides a \emph{Renderer} per playback.
  \item The renderer calls the emulation engine by blocks (about 100~ms), resamples to the output rate (48~kHz) and
        applies gain, fades and loop limits.
  \item The Android service writes the result to the audio output. Emulation state snapshots, taken regularly, make
        backward seeking fast without replaying from the start.
\end{enumerate}

\section{Visualisation data}
Every engine exposes the same three streams through the \emph{VoicesScope} interface:
\begin{description}
  \item[Layout] list of voice groups (one chip per group) and their names, sent on every change.
  \item[Per-voice outputs] samples of each voice at the output rate, resampled like the sound.
  \item[Voice states] snapshot of each voice about every millisecond: frequency, level, flags, chip kind and 64 bytes
        specific to the chip (registers, envelopes, wavetable\dots).
\end{description}
For libvgm chips these streams come from an extension added to the emulation cores: a per-sample callback (per-voice
outputs) and a state function. reSID and reSIDfp, ares and Ootake received equivalent extensions.

\section{Rendering}
\begin{itemize}
  \item A worker thread prepares each frame: for each voice, the time window is aligned on the fundamental period of
        the note (trigger), then reduced to two points per pixel column (minimum and maximum), so no peak is lost.
  \item Traces are drawn by Vulkan when available, otherwise OpenGL ES, with thin antialiased lines.
  \item The dashboard is drawn in two layers: a static layer (frames, captions), redrawn only when the layout
        changes, and a dynamic layer (values, curves), redrawn every frame.
\end{itemize}

\section{Tests}
\begin{itemize}
  \item More than 1~000 JVM tests (Robolectric) for the application: models, playlist storage, view computations, and
        render tests producing the screenshots of this manual.
  \item Native check tools: oscilloscope measurement (\texttt{scope\_bench}), per-voice outputs of VGM chips, SID
        parameters, seeking.
\end{itemize}
